\printedpage{54}
% head: A. Grothendieck : Dérivateurs, Ch. XV, Édition des Universités de Montpellier II et Paris VII
% foot: Version 5 novembre 2010
are satisfied. Necessity is obvious, by definition of closed Quillen triples, and for sufficiency, we note that it follows from theorem 7.1, 2$^\circ$, and from the part of 7.2 already proved, that $(C,TF)$, $(TC,F)$ are closed Quillen couples, and we have moreover
\[
TC \;=\; C\cap W\ , \qquad TF \;\subset\; F\cap W
\]
(since $TF\subset W$ by (7.2.7), and $TF\subset F$ since $TC\subset C$). We know
\msp{page 98}
(paragraph 4.3, corollary 4, page 36) that it follows that $(W,C,F)$ is a closed Quillen triple. This completes the proof of 7.1 and 7.2.

Let us sum up the essentials:

\textbf{Corollary 7.4.}\footnote[35]{\textbf{N.B.} It is ultimately this corollary that is the main result of the present section.} \textit{Let $\mathcal{M}$ be a category satisfying the conditions} (a), (b), (c), (d), (e) (p.~76) \textit{and} (e$'$), (f) ((7.2.1), p.~88), \textit{finally} (g) (7.2.6), i.e. \textit{stability of the set} Mon \textit{of monomorphisms of $\mathcal{M}$ under cobase change. Let $W\subset \Fl(\mathcal{M})$ be a localizer in $\mathcal{M}$, which is} stable under filtering $\varinjlim$ \textit{and moreover} (\textit{a subtler condition, but in practice innocuous}) accessible. \textit{We put}
\[
C \;=\; \Mon\ , \qquad TC \;=\; C\cap W\ , \qquad TF \;=\; C_*\ , \qquad F \;=\; (TC)_*\ .
\]
\textit{Then:}
\begin{itemize}[leftmargin=*]
\item[1$^\circ$)] \textit{$(C,TF)$ is a Quillen couple} (\textit{a statement independent of the datum $W$}), \textit{and even a Quillen couple with functorial factorization.}
\item[2$^\circ$)] \textit{For $(TC,F)$ to be a Quillen couple, it is necessary and sufficient that $TC$ be stable under cobase change.} (\textit{And then there is functorial factorization.})\footnote[36]{\textbf{N.B.} In cases 1$^\circ$ and 2$^\circ$, $C$, $TC=C\cap W$ are Q-generated on the left by \textit{small} subsets.}.
\item[3$^\circ$)] \textit{For $(W,C,F)$ to be a closed Quillen triple, it is necessary and sufficient that the following conditions be satisfied.}
\[
\left\{
\begin{aligned}
&\text{a)}\quad TC \overset{\text{def}}{=} W\cap \Mon\ \textit{stable under cobase change.}\\
&\text{b)}\quad TF \overset{\text{def}}{=} (\Mon)_* = \textit{maps of injective type} \subset W.
\end{aligned}
\right.
\tag*{(7.4.1)}
\]
\item[4$^\circ$)] \msp{page 99}\textit{For $(W,C,F)$ to be a closed and left proper Quillen triple} (i.e. $C\subset \Cof_W$), \textit{it is necessary and sufficient that we have condition} (7.4.1), (b) (i.e. \textit{$W$ contains the morphisms of injective type}), \textit{and moreover}
\[
\begin{aligned}
&C\subset \Cof_W,\ \textit{or again: for every cocartesian square}\\
&\qquad\qquad\begin{tikzcd}[ampersand replacement=\&, row sep=large, column sep=large]
A \arrow[r, "i"] \arrow[d, "\alpha"'] \& B \arrow[d, "\beta"] \\
A' \arrow[r, "i'"] \& B'\ ,
\end{tikzcd}\\
&\textit{with $i$ a monomorphism, then } \alpha\in W \Longrightarrow \beta\in W.
\end{aligned}
\tag*{(7.4.2)}
\]

\item[5$^\circ$)] \printedpage{55}
% head: A. Grothendieck : Dérivateurs, Ch. XV, Édition des Universités de Montpellier II et Paris VII
% foot: Version 5 novembre 2010
\textit{For $(W,C,F)$ to be a closed Quillen triple and} proper, \textit{it is necessary and sufficient that the two conditions of 4$^\circ$ be satisfied, and also that we have:}
\[
\begin{aligned}
&F\subset \Fib_W,\ \textit{or again: for every cartesian square}\\
&\qquad\qquad\begin{tikzcd}[ampersand replacement=\&, row sep=large, column sep=large]
X \arrow[r, "q"] \arrow[d, "f"'] \& X' \arrow[d, "f'"] \\
Y \arrow[r, "p"'] \& Y'\ ,
\end{tikzcd}\\
&\textit{if } f\in F \ (\text{i.e. } f\in (C\cap W)_*)\ \textit{and } p\in W,\ \textit{we have } q\in W.
\end{aligned}
\tag*{(7.4.3)}
\]
\end{itemize}

\textsc{Proof of} 7.4. Parts 1$^\circ$, 2$^\circ$, 3$^\circ$ are just a summary of what has been established. For 4$^\circ$, we have to see that the stated conditions imply (7.4.1), $(a)$, now we have $C\subset \Cof_W$, hence $C\cap W\subset \Cof_W\cap W = W^{\univ}$
\msp{page 100}
hence $C\cap W = C\cap W^{\univ}$, which is stable under cobase change, since $C$ and $W^{\univ}$ are. As for 5$^\circ$, it is a tautology, taking 4$^\circ$ into account.

\section*{8\quad Application to the categories of models $(A\chap,W_A)$}

Let $A$ be in Cat, and let us consider in $A\chap$ the localizer $W_A$, inverse image under
\[
\left\{
\begin{aligned}
i_A : A\chap \quad &\longrightarrow \quad \Cat\\
F \quad &\longmapsto \quad A/F
\end{aligned}
\right.
\tag*{(8.1)}
\]
of a given basic localizer $W$, satisfying the conditions
\[
\mathrm{W}\,(\underbrace{1}_{\text{sat.}},\underbrace{2}_{\text{loc.}},\underbrace{3}_{\text{final obj.}}\,,\underbrace{5}_{\text{conn. }\pi_0}\,,\underbrace{6}_{\varinjlim},\underbrace{7\text{ bis}}_{\substack{\text{coc. squares}\\ \text{or ``M.V.''}}}\,,\underbrace{8}_{\text{acc.}}),\quad
\text{\parbox[t]{4.3cm}{where W\,(8) is the accessibility of $W$ (thus $W$ is L-accessible if it also satisfies W\,(6)).}}
\tag*{(8.2)}
\]

To be able however to apply the developments on accessible subsets of a big category in §5, we must verify the conditions $(a)$, $(b)$, $(c)$, $(d)$, $(e)$ of §5, both in Cat, and (especially) in $A\chap$.

a) Stability under small filtering $\varinjlim$, O.K.

b) Cardinal filtration. Well known for $A\chap = \underline{\mathrm{Hom}}(A,\mathrm{Ens})$, starting from that of Ens by the cardinals (\textit{cf.} SGA 4 I, 9.24, p.~177). For Cat, we define a cardinal filtration
\[
(\Cat^{\pi})_{\pi\geq\pi_0}\ , \qquad \text{where } \pi_0 = \text{cardinal of countable sets}\ ,
\tag*{(8.3)}
\]
\printedpage{56}
% head: A. Grothendieck : Dérivateurs, Ch. XV, Édition des Universités de Montpellier II et Paris VII
% foot: Version 5 novembre 2010
\msp{page 101}
by the condition, for $A$ in Cat,
\[
A\in \mathrm{Filt}^{\pi}(\Cat) \;=\; \Cat^{\pi} \qquad\Longleftrightarrow\qquad \operatorname{card}\Fl(A) \;\leq\; \pi\ .
\tag*{(8.4)}
\]

We must verify the conditions $(a)$, $(b)$, $(c)$ of \textit{loc. cit.} 9.12. Condition $(a)$ is obvious, $(b)$ is contained in the stronger condition of stability of $\Cat^{\pi}$ under $\varinjlim\limits_{I}$, for $I$ a category such that $\operatorname{card}\Fl(I)\leq\pi$ (this is condition $(d_0)$ of these notes, p.~45-46). Let us note that $\Cat^{\pi}$ is also stable under finite $\varprojlim$ (and if $\pi$ is of the form $2^c$, it is even stable under inverse limits $\varprojlim\limits_{I}$, with $\operatorname{card}\Fl(I)\leq c\ \ldots$). Condition $(c)$ is easily verified by the method of §5, noting that for $X$ in Cat, the set
\[
\mathrm{Ss}_{\pi}(X)
\]
of subcategories of $X$ which are in $\Cat^{\pi}$ is upward filtering and large relative to $\pi$, and that $X$ is its direct limit. On the other hand, it is also clear that if $X\in\Cat^{\pi'}$ with $\pi'>\pi$, then
\[
\operatorname{card}\mathrm{Ss}_{\pi}(X) \;\leq\; \pi'^{\pi}\ ,
\]
which finishes establishing condition $(c)$.

Moreover, condition $(d)$ on p.~45 is immediately verified for the $\Cat^{\pi}$, and also for the $A^{\wedge\pi}$ (for $\pi\geq\operatorname{card}\Fl(A)$).

Also that monomorphisms are stable under filtering $\varinjlim$ (condition $(e)$, page 46).

\msp{page 102}
Let us pass to the additional conditions $(e')$, $(f)$, $(g)$ required in corollary 7.4. Condition $(e')$ (commutation of filtering $\varinjlim$ with base change) is satisfied in $A\chap$ and in Cat. Condition $(f)$ says that for two subobjects $A_1$, $A_2$ of an object $B$, if $A_0=A_1\cap A_2$, then
\[
A_1\amalg_{A_0}A_2 \longrightarrow B
\]
is a monomorphism. I doubt that this is true in Cat, and it does not matter, for it is to $A\chap$, not to Cat, that we shall want to apply 7.4 -- and it is obviously true in $A\chap$, and in any topos, being so in (Ens). Finally, stability of monomorphisms under cobase change is surely false in Cat, but true in $A\chap$ -- so it is O.K.

Since $W$ is stable under filtering $\varinjlim$ (W\,(6)), and $i_A$ commutes with them, $W_A$ is also stable. Similarly, $W_A$ is L-accessible, by the

\textbf{Lemma 8.1.} \textit{Let $\mathcal{M}$, $\mathcal{M}'$ be two categories satisfying} (a), (b), (c), (d), (e), \textit{and $\varphi:\mathcal{M}\longrightarrow\mathcal{M}'$ a functor commuting with filtering $\varinjlim$. Then for any L-accessible subset $M'$ of $\mathcal{M}'$, $\varphi^{-1}(M')$ is L-accessible.}

\msp{page 103}
\textit{To be precise, let $\pi$ be a cardinal such that $M'$ is $\pi$-accessible and $\varphi$ maps $\mathcal{M}^{\pi}$ into $\mathcal{M}'^{\pi}$ (for example, for this last condition, it suffices $\pi=2^c$ for $c$ large enough,} cf. loc. cit. 9.14, p.~153). \textit{Then $M=\varphi^{-1}(M')$ is $\pi$-accessible.}

Immediate verification.

We therefore have the

\printedpage{57}
% head: A. Grothendieck : Dérivateurs, Ch. XV, Édition des Universités de Montpellier II et Paris VII
% foot: Version 5 novembre 2010
\textbf{Lemma 8.2.} \textit{Let $A$ be in} Cat. \textit{Then $(A\chap,W_A)$ satisfies the preliminary conditions of corollary} 7.4, \textit{as well as the following conditions.}
\[
\underbrace{\Mon(A\chap)}_{\substack{\textit{set of}\\ \textit{monomorphisms}\\ \textit{in } A\chap}} \quad\subset\ \Cof_{W_A}\ ,
\tag*{(8.2.1)}
\]
\[
\Mon\cap W_A \qquad \textit{stable under cobase change}\ ,
\tag*{(8.2.2)}
\]
\textit{in other words} (\textit{taking into account that} Mon \textit{is stable under cobase change}), \textit{for a cocartesian square}
\[
\begin{tikzcd}[row sep=large, column sep=large]
A \arrow[r, hook, "i"] \arrow[dr, "\alpha"'] & B \arrow[dr, "\beta"] & \\
& A' \arrow[r, hook, "i'"] & B'\ ,
\end{tikzcd}
\tag*{(8.2.3)}
\]
\textit{with $i$ a monomorphism, if $\alpha\in W_A$ (resp. $i\in W_A$), we have $\beta\in W_A$ (resp. $i'\in W_A$). In fact, we have better:}
\[
i\ \in\ W_A \ \Longleftrightarrow\ i'\ \in\ W_A\ , \qquad \alpha\ \in\ W_A \ \Longleftrightarrow\ \beta\ \in\ W_A\ .
\tag*{(8.2.4)}
\]

\msp{page 104}
It remains to verify (8.2.4), which was seen in XII \ldots

We therefore conclude from 7.4 1$^\circ$, 2$^\circ$, 3$^\circ$:

\textbf{Corollary 8.3.} \textit{Let us put}
\[
\left\{
\begin{aligned}
C \;&=\; \Mon(A\chap)\ , & TC \;&=\; C\cap W\ ,\\
F \;&=\; (TC)_*\ , & TF \;&=\; C_*\ .
\end{aligned}
\right.
\tag*{(8.3.1)}
\]
\textit{Then $(C,TF)$ and $(TC,F)$ are Quillen couples with functorial factorization.} (\textbf{N.B.} \textit{$C$ and $TC$ being moreover Q-generated on the left by small subsets of} $\Fl(A\chap)$.) \textit{For $(W_A,C,F)$ to be a closed Quillen triple, it is necessary and sufficient that we have}
\[
\underbrace{TF}_{\overset{\text{def}}{=}\,(\Mon)_*} \quad\subset\ W_A\ ,
\tag*{(8.3.2)}
\]
i.e. \textit{that every arrow $f:X\longrightarrow Y$ of injective type in $A\chap$} (i.e. \textit{which makes $X$ an injective object of $A\chap/Y$}) \textit{be in $W_A$.}

\textit{When this condition is satisfied, the Quillen triple $(W_A,C,F)$ is left proper, by} 8.2.1.

It is not clear that it is necessarily proper-plain-and-simple, \textit{i.e.} that we have
\[
F\ \subset \Fib_W \qquad (?)\ .
\tag*{(8.3.3)}
\]

We will come back to this later. For the moment, the essential thing is to examine the condition (8.3.2) on $A$.

\printedpage{58}
% head: A. Grothendieck : Dérivateurs, Ch. XV, Édition des Universités de Montpellier II et Paris VII
% foot: Version 5 novembre 2010
\msp{page 105}
\textbf{Theorem 8.4.} \textit{Let $W$ be a basic localizer in} Cat \textit{satisfying} $\mathrm{W}\,(1,2,3,6,7,8)$, \textit{notations of} XII, (i.e. \textit{all the axioms considered so far, except perhaps the one} W\,(4) \textit{on fibrations} \texttt{[and the one W\,(5) on connectedness]}), \textit{let us suppose moreover $W$ stable under direct factors} (\textit{strengthening} W\,(1\,bis) \textit{of the axiom of mild saturation} W\,(1)).\footnote[37]{\textbf{N.B.} For the validity of 1$^\circ$, 3$^\circ$ we do not need W\,(6) ( $\varinjlim$ ) or W\,(8) (accessibility), nor even W\,(7) (cocartesian square axiom). The latter is already necessary to ensure part $(a)$ of corollary 8.5. But for part 2$^\circ$ of theorem 8.4, we need \textit{all} the axioms of the basic localizer $W$, except only W\,(4) \texttt{[and W\,(5)]}.} \textit{Let $A$ be an object of} Cat, \textit{and $W_A=i_A^{-1}(W)$ the localizer of $A\chap$ deduced from $W$. With the notations of} 8.3:
\begin{itemize}[leftmargin=*]
\item[1$^\circ$)] \textit{The following conditions are equivalent.}
\begin{itemize}[leftmargin=*]
\item[a)] \textit{We have} (8.3.2),
\[
TF\ \subset\ W_A\ ,
\]
i.e. \textit{every arrow of injective type in $A\chap$ is in $W_A$.}
\item[b)] \textit{Let}
\[
L \;=\; L_A
\tag*{(8.4.1)}
\]
\textit{be the Lawvere element in $A\chap$, hence the presheaf}
\[
a \quad\longmapsto\quad
\begin{aligned}
&\textit{set of subpresheaves of $a$}\\
&(\text{i.e. }\textit{set of opens of $A/a$},\\
&\text{i.e. }\textit{of sieves of $A/a$})\ ,
\end{aligned}
\]
\textit{and let us consider the arrow $L\longrightarrow e_A$ (final object of $A\chap$). We want this arrow to be $W_A$-aspheric,} i.e. \textit{to be universally in $W_A$ -- or what comes to the same thing, that for every $a\in A$, $L_A\times a$ $(=L_{A/a})$ be a $W_A$-aspheric object}
\msp{page 106}
\textit{of $A\chap$ (or, what comes to the same thing, of $(A/a)\chap$).}
\item[c)] \textit{There exists an object $I$ of $A\chap$ such that $I\longrightarrow e_A$ is $W_A$-aspheric, with two sections $\delta_0$, $\delta_1$ of $I$ over $e_A$, with} $\mathrm{Ker}(\delta_0,\delta_1)=\emptyset_A$ \textit{(empty presheaf).}
\item[d)] \textit{The localized categories $A/a=B$ are} test categories \textit{for $W$} (cf. \textit{Pursuing Stacks}, \textit{§65, especially pages 172-173), which amounts to saying that we have this, for each of the categories $B=A/a$:}
\[
\begin{aligned}
&\textit{The functors}\\
&\left\{
\begin{aligned}
&i_B:\ \ B\chap \ \longrightarrow\ \Cat\ \ (\ F\ \longmapsto\ B/F\ )\\
&j_B:\ \ \Cat\ \longrightarrow\ B\chap\ \ (\ C\ \longmapsto\ (b\longmapsto \mathrm{Hom}(B/b,C))\ )
\end{aligned}
\right.
\end{aligned}
\tag*{(8.4.2)}
\]
(\textit{where $j_B$ is the right adjoint of $i_B$}) \textit{are both ``model preserving''}, i.e.
\[
\left\{
\begin{aligned}
W_B \;&=\; i_B^{-1}(W) \qquad \textit{(by definition of $W_B$, in fact)}\\
W \;&=\; j_B^{-1}(W_B)\ ,
\end{aligned}
\right.
\tag*{(8.4.3)}
\]
\end{itemize}
\end{itemize}